\chapter[Vulnerability Analysis]{Vulnerability Analysis}
\label{cp:vuln-analysis}

\autoref{cp:system-design} established the directed dependency graph that SBOM Lens constructs from a CycloneDX SBOM. Every node in that graph carries a \texttt{purl} attribute --- a Package URL uniquely identifying the component --- as defined in \autoref{sec:node-representation}. \autoref{cp:vuln-analysis} uses those PURL attributes as the anchor for vulnerability matching: by querying the Open Source Vulnerabilities database (OSV) with the PURLs present in the graph, SBOM Lens enriches the structural model with real-world risk data. OSV aggregates advisories from the GitHub Advisory Database, the PyPA advisory database and many other upstream sources, so a single PURL-native query reaches all of them. The chapter is organised into five sections. \autoref{sec:multi-source} describes why OSV is the vulnerability source and how the integration pipeline works. \autoref{sec:dedup-storage} describes how vulnerability records are stored in the in-memory \texttt{VulnStore} and deduplicated by CVE ID and PURL pair. \autoref{sec:cvss-computation} explains how CVSS v3.1 base scores are computed on the client side and maps scores to the severity bands used throughout the system. \autoref{sec:vuln-workflow} covers the analyst-facing management workflow: automated fetching on SBOM upload and the three-state status lifecycle exposed through the REST API. \autoref{sec:graph-overlay} describes how vulnerability data is projected back onto the Cytoscape graph as three distinct overlay types --- severity colour coding, blast-radius highlighting, and dominator-tree highlighting --- giving analytical meaning to the CSS overlay classes introduced in \autoref{sec:cytoscape-rendering}.


\section{Vulnerability Source}
\label{sec:multi-source}

Matching SBOM components against known vulnerabilities requires querying an external database that maintains curated records of published vulnerabilities. This section justifies the choice of OSV as that database and describes the integration pipeline.


\subsection{Rationale for OSV}
\label{sec:multi-source-rationale}

No single vulnerability database covers all software ecosystems with equal completeness, and databases differ in how precisely they identify the affected package. NVD, maintained by NIST, identifies products with Common Platform Enumeration (CPE) strings, a scheme designed for commercial software. As noted in \autoref{sec:cve-nvd}, matching open-source packages against CPE strings is imprecise: a package may have several CPE strings or none, and name-based matching can return vulnerabilities for unrelated products that share a name. OSV, by contrast, was designed from the outset to accept PURL queries and to express affected versions as ranges, which makes it precise for exactly the identifiers an SBOM carries.

OSV is also an aggregator. It ingests the GitHub Advisory Database (GHSA), the PyPA advisory database, RustSec, the Go vulnerability database, and many other upstream sources, and each entry cross-references CVE and GHSA identifiers in its \texttt{aliases} field \citep{OSV}. A single OSV query therefore reaches the advisory data that separate GHSA or NVD queries would have provided for open-source packages, without the imprecise name-to-product mapping that NVD requires. Finally, OSV requires no API key, so SBOM Lens produces vulnerability results in any environment without prior registration.

SBOM Lens therefore uses OSV as its sole vulnerability source. The cost of this choice is dependence on OSV's aggregation: an advisory published upstream reaches SBOM Lens only once OSV has ingested it, and vulnerabilities that exist only in NVD, without an open-source advisory, are not reported. Direct integration of further sources is discussed as future work in \autoref{cp:conclusion}.


\subsection{OSV Integration Pipeline}
\label{sec:osv-pipeline}

The OSV integration is implemented as a two-phase pipeline in \texttt{fetchVulnerabilitiesForSBOM()} inside \texttt{vulnerabilityAPI.js}. The first phase issues a single batch request to the \texttt{/v1/querybatch} endpoint with an array of PURL queries --- one entry per component in the loaded SBOM. This batch design avoids \(N\) serial HTTP requests for an SBOM with \(N\) components: the entire component list is submitted in one round-trip, and the response returns the matching OSV vulnerability IDs for each PURL position.

The second phase resolves each unique OSV ID to its full advisory record by fetching \texttt{/v1/vulns/\{id\}}. The detail record contains the canonical summary, published and modified dates, a list of CVSS severity objects, and an \texttt{aliases} array that may include CVE identifiers. CVE alias resolution is performed by scanning the \texttt{aliases} array for an entry beginning with the prefix \texttt{"CVE-"}: the first such alias becomes the canonical \texttt{cveId} stored in \texttt{VulnStore}. If no CVE alias is present, the OSV identifier itself is used as the \texttt{cveId} fallback; deduplication logic (\autoref{sec:dedup-logic}) still operates correctly on OSV IDs.

Severity is derived in three steps. If the \texttt{severity[]} array of the detail record contains an entry whose \texttt{type} is \texttt{"CVSS\_V3"}, its \texttt{score} field --- a CVSS vector string rather than a pre-computed number --- is passed to \texttt{cvssV3BaseScore()} (\autoref{sec:cvss-implementation}) and the resulting base score is mapped to a severity band. CVSS v4 vectors are not scored, because the v3.1 formula does not apply to them. If no v3 vector is present, the severity label that the originating advisory database attaches to the record (\texttt{database\_specific.severity}, as published by GHSA) is used, with GHSA's \textsc{moderate} mapped to \textsc{medium}. If neither is available, the record's severity is \textsc{unknown}; it is never silently treated as \textsc{low}.

The pipeline is invoked automatically through the \texttt{useVulnerabilityFetcher} React hook whenever an SBOM is uploaded to the library. The hook follows a fire-and-forget pattern: it is asynchronous, no error is ever shown to the analyst, and the graph view is never blocked waiting for vulnerability data to arrive. Records are injected into \texttt{VulnStore} sequentially rather than in parallel: submitting all records in parallel would cause duplicate POSTs to race the backend deduplication check before it has processed the first record. When an upload contains pairs already in the store --- because another SBOM shares the component, or because the same SBOM is uploaded again after being deleted --- the backend rejects each with HTTP 409, and the frontend hook treats 409 as a no-op, so the analyst observes no duplication.


\section{Deduplication and Storage}
\label{sec:dedup-storage}

The same vulnerability can reach the store more than once: the fetch runs for every SBOM uploaded, and SBOMs in one library often share components; and an analyst can also enter records manually. \texttt{VulnStore} is the in-memory store that absorbs these records and ensures that each unique vulnerability-component pair is stored exactly once. This section describes the storage schema and the deduplication logic.


\subsection{VulnStore Design}
\label{sec:vulnstore-design}

\texttt{VulnStore} is implemented as a Python class in \texttt{backend/store/vuln\_store.py}. It maintains a single in-memory dictionary, \texttt{self.vulns}, keyed by UUID strings generated at record creation time with \texttt{str(uuid.uuid4())}. State is ephemeral: the dictionary is cleared on server restart, consistent with the ephemeral design of \texttt{SBOMStore} described in \autoref{sec:in-memory-store} and appropriate for a research prototype evaluated on a fixed dataset.

Each record in \texttt{self.vulns} conforms to a consistent schema:

\begin{itemize}
    \item \textbf{id} --- UUID string uniquely identifying the record.
    \item \textbf{purl} --- Package URL of the affected component, matching the node attribute described in \autoref{sec:node-representation}.
    \item \textbf{cveId} --- CVE identifier or OSV identifier if no CVE alias was resolved.
    \item \textbf{severity} --- severity band (\textsc{critical}, \textsc{high}, \textsc{medium}, \textsc{low}, or \textsc{unknown}) derived as described in \autoref{sec:osv-pipeline}.
    \item \textbf{status} --- current remediation status; one of \texttt{NOT\_STARTED}, \texttt{IN\_PROGRESS}, or \texttt{REMEDIATED}.
    \item \textbf{patchAvailable} --- boolean, set by the analyst when entering a record manually, indicating that a patched version is known to exist; records from the OSV fetch always carry \texttt{false}.
    \item \textbf{notes} --- free-text analyst annotation.
    \item \textbf{source} --- origin of the record: \texttt{"osv"} for the automatic fetch or \texttt{"Manual"} for analyst entry.
    \item \textbf{publishedDate} --- date the advisory was first published.
    \item \textbf{lastUpdated} --- server-side timestamp refreshed on each \texttt{PUT} update.
    \item \textbf{timeline} --- list of status-transition objects, each carrying a timestamp, the new status value and a note, recording the remediation history of the record.
\end{itemize}

The \texttt{source} field preserves the provenance of each record, allowing an analyst to distinguish vulnerabilities found by the automatic OSV scan from those entered manually.


\subsection{Deduplication Logic}
\label{sec:dedup-logic}

Deduplication is enforced in the \texttt{VulnStore.add()} method before any new record is committed to the store. The dedup key is the pair of \texttt{cveId} and \texttt{purl}: the method searches the existing records for any entry where both fields match the incoming record.

\begin{verbatim}
duplicate = next((v for v in self.vulns.values()
    if v.get("cveId") == cve_id and v.get("purl") == purl), None)
\end{verbatim}

If a matching record is found, \texttt{add()} returns \texttt{None} without modifying the store. The REST endpoint wrapping this method responds with HTTP 409 Conflict. The \texttt{useVulnerabilityFetcher} hook treats a 409 response as a silent no-op: no error is raised, no duplicate appears in the store, and the analyst sees no change.

OSV IDs that could not be resolved to a CVE alias retain their OSV identifier as the \texttt{cveId} field; the deduplication logic operates on that OSV ID equally well, preventing duplicates when the same OSV entry is fetched again for another SBOM. Because the key is the first CVE alias, two distinct advisories that share that alias for the same component are also collapsed into one record, and the advisory stored second is discarded.

The first-written record for a pair is kept: a later record for the same pair, for example a manual entry with a different severity, is rejected rather than merged.


\section{CVSS v3.1 Base Score Computation}
\label{sec:cvss-computation}

Vulnerability records from the OSV pipeline carry CVSS vector strings rather than pre-computed numeric scores. SBOM Lens computes base scores from these vectors entirely on the client side, with no dependency on an external CVSS scoring service. This section describes how the CVSS v3.1 formula is applied, illustrates the computation with the Log4Shell vector, and explains the design of the custom JavaScript implementation.


\subsection{Formula and Metric Groups}
\label{sec:cvss-formula}

The CVSS v3.1 base score formula and the numerical values for each metric are fully defined in \autoref{sec:cvss}; this section does not repeat that derivation, but states explicitly the severity band thresholds that SBOM Lens uses to map numeric scores to visual encodings.

SBOM Lens maps CVSS base scores to four severity bands, plus \textsc{unknown} for records without any usable severity (\autoref{sec:osv-pipeline}):

\begin{center}
    \begin{tabular}{ll}
        \toprule
        \textbf{Severity} & \textbf{Score Range} \\
        \midrule
        \textsc{critical} & \(\geq 9.0\) \\
        \textsc{high} & \(7.0\)--\(8.9\) \\
        \textsc{medium} & \(4.0\)--\(6.9\) \\
        \textsc{low} & \(< 4.0\) \\
        \bottomrule
    \end{tabular}
\end{center}

These thresholds follow the NVD/NIST severity rating scale. FIRST.org's published five-band scheme adds a ``None'' band for a base score of exactly \(0.0\); SBOM Lens collapses this band into \textsc{low} to simplify the overlay colour mapping to four distinct visual classes. When the input to \texttt{parseCVSSSeverity()} is missing or cannot be scored, the function returns \textsc{unknown} rather than defaulting to \textsc{low}, so that missing severity data is not mistaken for low risk. On the graph, \textsc{unknown} nodes are drawn with the \textsc{low} colour; in the data, and in the evaluation of \autoref{cp:evaluation}, they remain a separate band.

The scope-sensitivity of the Privileges Required (PR) metric is handled explicitly: when Scope is Changed (\texttt{S:C}), the PR Low coefficient is \(0.68\) rather than \(0.62\), and the PR High coefficient is \(0.50\) rather than \(0.27\); PR None is \(0.85\) in both cases. The custom implementation branches on the parsed Scope value when looking up PR coefficients, matching the FIRST.org specification precisely.


\subsection{Worked Example: Log4Shell}
\label{sec:cvss-log4shell}

Log4Shell (\texttt{CVE-2021-44228}) was introduced as the principal illustrative case study in \autoref{sec:cvss} and analysed from a supply-chain perspective in \autoref{sec:case-log4shell}. \autoref{cp:vuln-analysis} applies the same vulnerability to the implementation: given its CVSS vector, the client-side scoring function must produce the expected base score.

The Log4Shell vector is:
\begin{center}
    \texttt{CVSS:3.1/AV:N/AC:L/PR:N/UI:N/S:C/C:H/I:H/A:H}
\end{center}

Substituting the numerical values specified in \autoref{sec:cvss} gives: \(AV = 0.85\), \(AC = 0.77\), \(PR = 0.85\) (Scope Changed branch), \(UI = 0.85\), \(C = 0.56\), \(I = 0.56\), \(A = 0.56\). The Impact Sub-Score (ISS) is:
\[
    \mathit{ISS} = 1 - (1 - 0.56)(1 - 0.56)(1 - 0.56) = 1 - 0.44^3 = 1 - 0.0852 \approx 0.9148
\]

Because Scope is Changed, the Impact term uses the modified formula:
\begin{align*}
    \mathit{Impact} &= 7.52 \times (\mathit{ISS} - 0.029) - 3.25 \times (\mathit{ISS} - 0.02)^{15} \\
    &\approx 7.52 \times 0.8858 - 3.25 \times 0.8948^{15} \\
    &\approx 6.661 - 0.614 \approx 6.048
\end{align*}

The Exploitability term is:
\[
    \mathit{Exploitability} = 8.22 \times 0.85 \times 0.77 \times 0.85 \times 0.85 \approx 3.887
\]

The base score is therefore:
\begin{align*}
    \mathit{Base} &= \min\bigl(1.08 \times (6.048 + 3.887),\ 10\bigr) \\
    &= \min(1.08 \times 9.935,\ 10) \\
    &= \min(10.73,\ 10) = 10.0
\end{align*}

The CVSS roundup rule produces a final score of \(10.0\), which falls in the \textsc{critical} band (\(\geq 9.0\)). In the graph overlay described in \autoref{sec:severity-colour}, a Log4Shell-affected component node is therefore assigned the \texttt{.vuln-critical} CSS class and rendered in red.


\subsection{Custom Client-Side Implementation}
\label{sec:cvss-implementation}

The CVSS scoring logic is implemented entirely in \texttt{vulnerabilityAPI.js} with no external CVSS library dependency. This is a deliberate design choice: adding a CVSS computation library to the frontend \texttt{package.json} would introduce a runtime dependency whose licensing, maintenance status, and specification compliance would need independent verification. Implementing the formula from the FIRST.org specification in approximately seventy lines of JavaScript keeps the frontend self-contained and makes the scoring logic directly auditable alongside the rest of the source code.

The function \texttt{cvssV3BaseScore(vector)} accepts a CVSS vector string, splits it on the \texttt{"/"} separator to extract metric abbreviation--value pairs, and looks each pair up in a coefficient table to produce the numeric values used in the formulas above. It branches on the parsed Scope value to select the correct Impact and PR coefficients, computes the ISS, Impact, Exploitability, and Base sub-terms in sequence, and applies the CVSS roundup rule (\texttt{Math.ceil(base * 10) / 10}) before returning the final score. This roundup rule is specified by FIRST.org as ``the smallest number, specified to one decimal place, that is equal to or higher than its input''; the expression \texttt{Math.ceil(x * 10) / 10} implements it exactly.

The entry point for callers is \texttt{parseCVSSSeverity(scoreString)}, which accepts either a numeric score string (e.g.\ \texttt{"9.8"}) or a CVSS v3 vector string (as returned by OSV). If the input does not parse as a number and begins with \texttt{"CVSS:"}, it is treated as a vector and passed to \texttt{cvssV3BaseScore()}. This single entry point means the rest of the codebase does not need to distinguish numeric scores from vectors.


\section{Vulnerability Management Workflow}
\label{sec:vuln-workflow}

Beyond recording which components carry known CVEs, SBOM Lens provides a lightweight remediation-tracking workflow. Analysts can update the status of each vulnerability record, add notes, and observe how the record's status has changed over time. This section describes the automated fetch that populates the store on SBOM upload and the REST API that exposes the status lifecycle.


\subsection{Automated Fetch on SBOM Upload}
\label{sec:automated-fetch}

Vulnerability fetching is triggered automatically when an SBOM is uploaded to the library, without requiring any analyst action. The \texttt{useVulnerabilityFetcher} React hook receives the \texttt{components[]} array from the uploaded SBOM and passes the PURL list to the OSV pipeline described in \autoref{sec:osv-pipeline}.

The hook follows the fire-and-forget pattern introduced in \autoref{sec:osv-pipeline}: it is asynchronous, does not block rendering, and never surfaces an error to the analyst: a failed OSV request yields an empty result, and a backend error while storing records is logged with \texttt{console.warn} and abandons the remaining records of that SBOM. A network failure therefore never prevents the analyst from viewing the dependency graph, but an empty vulnerability list cannot be told apart from a failed lookup. Records are submitted to the backend's \texttt{POST /vulnerabilities} endpoint one at a time rather than in a batch request; this sequential submission prevents a parallel flood of identical records from racing the deduplication check on the server side.

Because records are keyed by CVE ID and PURL rather than by SBOM, an upload that shares components with an SBOM already in the library resubmits pairs that already exist in \texttt{VulnStore}; so does an SBOM uploaded again after being deleted, since deleting an SBOM does not delete its vulnerability records. Each such pair receives a 409 Conflict response, which the hook treats as a no-op. Vulnerabilities are fetched once, at upload: selecting an SBOM or reloading the page does not query OSV again, so an advisory published later reaches an SBOM only if it is uploaded again.


\subsection{Status Lifecycle and REST API}
\label{sec:status-lifecycle}

\texttt{VulnStore} exposes a full CRUD REST API through the \texttt{vulnerabilities} FastAPI router. The endpoints are:

\begin{itemize}
    \item \texttt{GET /vulnerabilities} --- returns all records in the store as a JSON array; no SBOM-scoping filter is applied, making the store a global view of all known vulnerabilities across all loaded SBOMs.
    \item \texttt{POST /vulnerabilities} --- creates a new record; returns HTTP 201 on success or HTTP 409 if the CVE ID + PURL pair already exists.
    \item \texttt{PUT /vulnerabilities/\{id\}} --- partial update of an existing record: every field sent replaces the stored value; the interface uses it to change \texttt{status} and extend \texttt{timeline}. The \texttt{lastUpdated} timestamp is set by the server.
    \item \texttt{DELETE /vulnerabilities/\{id\}} --- removes a single record.
    \item \texttt{DELETE /vulnerabilities} --- clears the entire store; useful for resetting state between evaluation runs without restarting the server.
\end{itemize}

The \texttt{status} field supports a three-state remediation lifecycle. \texttt{NOT\_STARTED} is the initial state assigned to every newly created record. \texttt{IN\_PROGRESS} indicates that a remediation effort is under way. \texttt{REMEDIATED} is the terminal state, indicating that the vulnerability has been patched or mitigated. The backend stores the field without validating it, so triage states such as a deliberate decision not to fix, or a false positive, could be added to the interface without changing the API; the evaluated version does not offer them. The first \texttt{timeline} entry is written by the server when the record is created; each later status change made through the interface appends an entry with a client-side timestamp and sends the extended list back. The server does not enforce this, since \texttt{PUT} replaces any field it receives, so the timeline is a remediation history for a single-user prototype rather than a tamper-proof audit trail. \autoref{fig:vulnerabilities-page} shows the Vulnerabilities page, where the analyst reviews the stored records and changes their status.

\begin{figure}[!htpb]
    \centering
    \screenshot{Figures/Chapter5/vulnerabilities-page}{Vulnerabilities page with the \texttt{code-server} SBOM selected: severity totals and the records table, with the Component column now filled in. Light theme, same window size as the other screenshots.}
    \caption[The Vulnerabilities page]{The Vulnerabilities page for the \texttt{code-server} SBOM: totals by severity band, and the stored vulnerability records with their component, severity band and remediation status.}
    \label{fig:vulnerabilities-page}
\end{figure}


\section{Graph Overlay Visualisation}
\label{sec:graph-overlay}

The vulnerability data held in \texttt{VulnStore} is projected back onto the Cytoscape.js graph through three distinct overlay types. Each overlay assigns one of the CSS classes pre-defined in \autoref{sec:cytoscape-rendering} --- \texttt{.vuln-critical}, \texttt{.vuln-high}, \texttt{.vuln-medium}, \texttt{.vuln-low}, \texttt{.blast-radius}, and \texttt{.dominator-tree} --- to the appropriate nodes, giving those classes the analytical meaning that \autoref{cp:system-design} deferred to this chapter. All three overlays are activated interactively by the analyst: the severity colour-coding overlay from the Analysis menu, and the blast-radius and dominator-tree overlays from the same menu, for the node chosen with \emph{Analyze Impact} in its context menu.


\subsection{Severity Colour Coding}
\label{sec:severity-colour}

The severity colour-coding overlay is activated from the Analysis menu once vulnerability records for the loaded SBOM are present in \texttt{VulnStore}. When it is active, each graph node is assigned a CSS overlay class determined by the highest-severity CVE that affects the component identified by that node's PURL.

A component may carry multiple CVE records in \texttt{VulnStore}; in that case the highest severity wins. A node whose PURL is associated with one \textsc{critical} and three \textsc{medium} vulnerabilities is displayed as \textsc{critical} --- the most dangerous exposure takes visual precedence. The Cytoscape overlay is applied programmatically: for each node, the vulnerability records whose \texttt{purl} field matches the node's PURL are retrieved from the store, the maximum of their severity bands (assigned by \texttt{parseCVSSSeverity()} when the records were fetched) is taken, and the node is assigned the matching class (\texttt{.vuln-critical}, \texttt{.vuln-high}, \texttt{.vuln-medium}, or \texttt{.vuln-low}). Nodes whose PURL matches no stored vulnerability record receive no severity class and keep their default colour. \autoref{fig:node-details} in \autoref{cp:system-design} shows the overlay applied to the \texttt{code-server} SBOM, with \texttt{basic-ftp}, its \textsc{critical} component, selected.



\subsection{Blast Radius Analysis}
\label{sec:blast-radius}

The blast-radius overlay is an interactive analysis triggered when the analyst chooses a node with \emph{Analyze Impact} in its context menu and switches the overlay on in the Analysis menu, as described in \autoref{sec:interaction-model}. The query is issued to \texttt{GET /\{sbom\_id\}/impact?node\_id=\{id\}}, which returns the blast radius, dependency paths, dominated nodes, and the HITS hub and authority scores of the queried node in a single response.

On the backend, the blast radius is computed as \texttt{list(nx.ancestors(G, node\_id)) + [node\_id]}. The semantic interpretation follows directly from the edge direction established in \autoref{sec:edge-representation}: an edge from A to B means A depends on B, so \texttt{nx.ancestors(G, B)} returns all nodes from which there exists a directed path to B --- the upstream consumers of B whose security posture is degraded if B is compromised; the blast radius is B together with these consumers. This framing is intentional and distinct from a ``downstream'' interpretation: in the SBOM Lens graph, the upstream consumers are the packages that pull in the vulnerable component, and so the packages that may be affected if the vulnerable node is compromised. Because the analysis is at component level, this over-approximates exposure: a consumer may never execute the vulnerable code (\autoref{sec:rw-call-graphs}).

The API response also includes up to five shortest paths from the root packages (components with in-degree zero) to the vulnerable node, computed with \texttt{nx.all\_shortest\_paths} from each root. These paths provide concrete dependency chains --- a sequence of ``package X uses package Y which uses the vulnerable package Z'' --- that an analyst can use to understand exactly how the vulnerable component entered the dependency tree. The interface lists them in the \emph{Dependency paths} section of the Details panel, together with the total number of paths from the root to the node and the number of consumers they pass through; selecting a path highlights it on the graph, and \emph{Highlight all paths} highlights every edge of the blast radius.

On the frontend, blast-radius nodes receive the \texttt{.blast-radius} CSS class (purple, as pre-defined in \autoref{sec:cytoscape-rendering}). \autoref{fig:blast-radius} shows the overlay for \texttt{basic-ftp}~5.0.3, the component with a \textsc{critical} advisory (CVE-2026-27699) in the \texttt{code-server} SBOM.

\begin{figure}[!htpb]
    \centering
    \screenshot[0.8\textwidth]{Figures/Chapter5/blast-radius}{SBOM Graph page, \texttt{code-server} SBOM: \emph{Analyze Impact} on \texttt{basic-ftp}, Blast Radius and severity overlays active, all paths highlighted. Light theme, same window size as the other screenshots.}
    \caption[Blast-radius overlay]{Blast radius of \texttt{basic-ftp} in the \texttt{code-server} SBOM. Its four consumers --- \texttt{get-uri}, \texttt{pac-proxy-agent}, \texttt{proxy-agent} and the application \texttt{code-server} --- are purple, \texttt{basic-ftp} keeps the red of its \textsc{critical} band, and the edges of the single dependency path are highlighted in orange. The Details panel lists the path, \texttt{code-server} \(\to\) \texttt{proxy-agent} \(\to\) \texttt{pac-proxy-agent} \(\to\) \texttt{get-uri} \(\to\) \texttt{basic-ftp}.}
    \label{fig:blast-radius}
\end{figure}

The impact API response also includes \texttt{hitsHub} and \texttt{hitsAuthority} scores of the analysed component; scores for every node are returned by \texttt{GET /sboms/\{id\}/global-analysis}. These scores quantify the component's structural centrality in the dependency network and are relevant to prioritising remediation effort; however, the HITS algorithm that produces them is described in detail in \autoref{cp:network-sci} (specifically \autoref{sec:hits}) and is not explained further here.


\subsection{Dominator Tree Analysis}
\label{sec:dominator-tree}

The dominator-tree overlay is returned in the same \texttt{GET /\{sbom\_id\}/impact} response as the blast radius, adding structural risk information without requiring an additional round-trip.

On the backend, dominance is computed from the root of the dependency graph, not from the vulnerable node. \texttt{nx.immediate\_dominators(G, root)} returns, for every component reachable from the root, its immediate dominator; when the SBOM has several root components, they are first joined under a virtual root so that a single dominator tree exists. The dominated set of the vulnerable node is then its entire subtree in that dominator tree: every component for which each path from the root passes through the vulnerable node. In supply-chain terms, a node with a non-empty dominated set is a structural chokepoint: if it were removed or replaced, all the dominated components would lose their connection to the application with no alternative route.

Dominated nodes are assigned the \texttt{.dominator-tree} CSS class on the frontend, displayed in cyan as pre-defined in \autoref{sec:cytoscape-rendering}. When both overlays are switched on for the same node, the analyst can read them together: nodes highlighted in purple are the upstream consumers affected by the vulnerability's propagation, while nodes highlighted in cyan are structurally dependent on the vulnerable node and would be isolated if it were removed. This dual encoding gives both a security-propagation view and an architectural-risk view in a single inspection, providing complementary information that neither overlay conveys alone. In the evaluated version, choosing a new node while both overlays are active refreshes only the blast radius, so the dominator-tree overlay must be switched off and on again to follow the new node. \autoref{fig:dominated-set} shows the dominated set of \texttt{express} in the \texttt{code-server} SBOM.

\begin{figure}[!htpb]
    \centering
    \screenshot[0.8\textwidth]{Figures/Chapter5/dominated-set}{SBOM Graph page, \texttt{code-server} SBOM: \emph{Analyze Impact} on \texttt{express}, Dominator Tree and severity overlays active. Light theme, same window size as the other screenshots.}
    \caption[Dominated set]{Dominated set of \texttt{express} in the \texttt{code-server} SBOM (cyan): the components reachable from the root only through \texttt{express}, which would lose their connection to the application if it were removed. The severity colours of the vulnerable components are also shown.}
    \label{fig:dominated-set}
\end{figure}
